Kretschmann's observation about covariance: its mathematical and physical
meaning Kretschmann's observation is one of the most important conceptual
points in the foundations of general relativity. It exposes a distinction that
is often overlooked: general covariance is not, by itself, a physical
principle. The central point is that virtually any physical theory can be
formulated in generally covariant form. Consequently, the fact that Einstein's
equations are generally covariant cannot, on its own, distinguish general
relativity from Newtonian mechanics, special relativity, or another physical
theory. This observation is associated with the German physicist Erich
Kretschmann, particularly his 1917 paper: Über den physikalischen Sinn der
Relativitätspostulate. Let us examine what this means, especially from the
perspective of covariance, invariance, and symmetry. 1. The original issue:
what is special about general covariance? Einstein's general theory of
relativity is formulated using tensor equations. For example, \[
    G_{\mu\nu}=\frac{8\pi G}{c^4}T_{\mu\nu}. \] This equation retains its form
under arbitrary smooth coordinate transformations: \[ x^\mu\longrightarrow
x'^\mu(x). \] Einstein regarded general covariance as an important
characteristic of the gravitational theory. The conceptual motivation was that
no coordinate system should have a privileged physical status. In particular,
the laws of physics should not depend on the arbitrary choice of coordinates.
But Kretschmann asked a very sharp question: If general covariance is such an
important physical criterion, why can't we formulate any physical theory in
generally covariant form? His answer was essentially that we can. 2. An
explicit example: special relativity in generally covariant form Consider
special relativity in its familiar formulation: \[ ds^2=c^2dt^2-dx^2-dy^2-dz^2.
\] The metric is the Minkowski metric, \[
    \eta_{\mu\nu}=\operatorname{diag}(1,-1,-1,-1). \] The equation of motion of
a free particle, using an affine parameter \(\lambda\), is \[
    \frac{d^2x^\mu}{d\lambda^2}=0. \] This formulation uses inertial Cartesian
coordinates. Now introduce arbitrary coordinates \(x^\mu\). The metric becomes
\[ ds^2=g_{\mu\nu}(x)dx^\mu dx^\nu. \] The free-particle equation becomes \[
    \boxed{ \frac{d^2x^\mu}{d\lambda^2} +\Gamma^\mu_{\alpha\beta}
    \frac{dx^\alpha}{d\lambda} \frac{dx^\beta}{d\lambda}=0. } \] This is the
geodesic equation. We can now write the equations of special relativity in a
form that is covariant under arbitrary coordinate transformations. We must,
however, retain the additional condition \[ \boxed{R^\mu{}_{\nu\alpha\beta}=0.}
\] This condition tells us that the metric is flat. Thus, we have obtained a
generally covariant formulation of special relativity: \[ \begin{cases}
    R^\mu{}_{\nu\alpha\beta}=0,\\[2mm] \nabla_\mu T^{\mu\nu}=0,\\[2mm]
\text{matter equations formulated using }g_{\mu\nu}. \end{cases} \] Nothing
physical has changed. We have merely changed the mathematical formulation. The
metric is still a prescribed, nondynamical background structure. It is not
determined by the matter distribution. This is precisely the kind of example
that illustrates Kretschmann's observation. sites.pitt.edu+1 3. Why does this
not make special relativity equivalent to general relativity? Here we reach the
essential point. Both theories can be expressed using arbitrary coordinates,
but their physical structures are different. Consider: Special relativity
General relativity Minkowski metric	Dynamical spacetime metric
\(R^\mu{}_{\nu\alpha\beta}=0\)	Curvature generally nonzero Metric prescribed
Metric determined by field equations Matter evolves on a fixed geometry	Matter
and geometry interact Poincaré spacetime symmetry	No fixed background metric
with a prescribed isometry group In general relativity, \[ G_{\mu\nu}[g]=\kappa
T_{\mu\nu} \] determines the metric dynamically. In special relativity, by
contrast, the metric is prescribed: \[ g_{\mu\nu}=\eta_{\mu\nu}. \] Even when
we express the latter in arbitrary coordinates, we have not made the metric
dynamical. This leads to a particularly important distinction: Coordinate
covariance is not the same as physical symmetry. The full diffeomorphism group
can be used to describe both theories, but the physical spacetime symmetries of
special relativity are the Poincaré transformations, not arbitrary
diffeomorphisms. The latter merely express the freedom to choose coordinates.
4. Passive covariance versus active covariance This distinction is especially
important. There are two different operations that are frequently confused. 4.1
Passive transformation We change the coordinate labels assigned to spacetime
events: \[ x^\mu\longrightarrow x'^\mu(x). \] The physical event remains the
same. Tensor components change according to the usual transformation laws. For
example, \[ g'_{\mu\nu}(x') = \frac{\partial x^\alpha}{\partial x'^\mu}
\frac{\partial x^\beta}{\partial x'^\nu} g_{\alpha\beta}(x). \] This is simply
a change of description. There is no new physical configuration. 4.2 Active
transformation Now consider a diffeomorphism \[ \phi:M\longrightarrow M. \]
Instead of changing coordinates, we transform the fields themselves. For
example, \[ g_{\mu\nu}\longrightarrow \phi^*g_{\mu\nu}. \] The transformed
metric generally represents a different assignment of geometrical fields to the
manifold's points. In general relativity, however, if we transform all
dynamical fields together, the transformed configuration is another solution of
the same field equations. The question then becomes: Are these two solutions
physically different? The modern answer is that solutions related by a
diffeomorphism represent the same physical geometry, provided we are treating
the diffeomorphism as a gauge transformation. This is a much stronger statement
than passive covariance. It is a statement about the physical interpretation of
the mathematical configurations. Norton's analysis is particularly useful here
because it explains how Kretschmann's objection and the physical importance of
diffeomorphism gauge freedom can both be correct. sites.pitt.edu+2 5. The
deeper issue: covariance versus symmetry Let us connect this with the
distinction between covariance and invariance. Suppose a theory is described by
equations \[ E[\Psi]=0, \] where \(\Psi\) denotes all the fields of the theory.
A transformation \(T\) is a symmetry of the theory if it maps solutions into
solutions: \[ E[\Psi]=0 \quad\Longrightarrow\quad E[T\Psi]=0. \] But we must be
careful: this condition alone does not distinguish physical symmetries from
gauge redundancies. For a global physical symmetry, the transformation relates
configurations that can represent physically different states, while preserving
the laws. For a gauge symmetry, the transformation relates different
mathematical descriptions of the same physical state. Now consider general
covariance. A generally covariant formulation means that the equations can be
expressed in tensorial form, retaining their form under arbitrary coordinate
transformations. But this does not necessarily imply that arbitrary
diffeomorphisms are physical symmetries. In particular, consider a theory with
a nondynamical background tensor \(B\). Its equations might be written
schematically as \[ E[\Psi;B]=0. \] Under an active diffeomorphism, we have \[
    \Psi\longrightarrow\phi^*\Psi, \qquad B\longrightarrow\phi^*B. \] The
equations are covariant if \[ E[\phi^*\Psi;\phi^*B] = \phi^*E[\Psi;B]. \]
However, if \(B\) is fixed, we must ask whether \[ E[\phi^*\Psi;B]=0. \] This
is generally not true. It is true only for transformations preserving the
background: \[ \phi^*B=B. \] These transformations form the symmetry group of
the background structure. For the Minkowski metric, this is the Poincaré group
(including the appropriate discrete transformations if the full group is
considered). This is an extremely useful way of understanding Kretschmann's
observation: - Covariance concerns the form of the equations under changes of
representation. - Symmetry concerns transformations preserving the physical
structure of the theory. - Gauge symmetry concerns transformations that leave
the physical state unchanged. These are three different notions. 6. Background
structures and the physical content of general relativity This brings us to a
more substantial interpretation of general relativity. A useful way to
characterize the difference between general relativity and special relativity
is through their background structures. In special relativity, the metric is
absolute in the sense that it is prescribed independently of the matter
dynamics. In general relativity, the metric is dynamical. This suggests that
the physically significant distinction is not simply general covariance, but
rather the status of the geometrical structures appearing in the theory. We can
schematically distinguish: \[ \begin{aligned} \text{SR:}&\quad
\eta_{\mu\nu}\text{ is prescribed},\\ \text{GR:}&\quad g_{\mu\nu}\text{ is
dynamical}. \end{aligned} \] This is closely related to the notion of absolute
objects, developed in the work of Anderson and Friedman. An absolute object is,
roughly speaking, a geometrical structure that is not dynamically determined by
the physical fields of the theory. There are technical subtleties in giving
this definition precisely, but the distinction is important. It also explains
why simply rewriting special relativity using arbitrary coordinates does not
turn it into general relativity. The Minkowski metric remains an absolute
background structure, even though its components are now coordinate-dependent.
ScienceDirect+1 7. What was Einstein's response? Einstein accepted the central
mathematical point of Kretschmann's objection. He nevertheless argued that
general covariance retained heuristic significance. The argument was
essentially that a generally covariant formulation could reveal physical
content previously hidden in the choice of preferred coordinates. A theory
might be expressible in arbitrary coordinates, but doing so could require
introducing additional geometrical structures and making its equations
considerably more complicated. The search for a generally covariant formulation
could therefore guide the discovery of a more natural theory. This is an
important qualification. General covariance may have little discriminatory
power as a purely formal requirement, but it can be a useful methodological
principle. John Norton discusses this distinction in detail in his historical
analysis of the controversy. IOPscience+1 8. Connection with the hole argument
There is another important consequence. Suppose we have a solution of
Einstein's equations, \[ (M,g_{\mu\nu},T_{\mu\nu}). \] Consider a
diffeomorphism \(\phi\) which differs from the identity only inside some
bounded spacetime region \(H\), called the hole. Then \[
    (M,\phi^*g_{\mu\nu},\phi^*T_{\mu\nu}) \] is also a solution. Outside the
hole, the two solutions coincide. Inside, their field components differ. If we
regard manifold points as possessing independent physical identities, we appear
to have two different physical evolutions compatible with identical initial
data. This was the source of Einstein's famous hole argument. The resolution is
to recognize that the two configurations are physically equivalent: the
manifold points have no independent physical identities apart from the
geometrical and material relations defined by the fields. In other words,
diffeomorphism invariance is associated with a redundancy in the mathematical
description. This is where the apparently negative Kretschmann observation
acquires a much more interesting modern significance. General covariance alone
does not distinguish GR from other theories. But the gauge freedom associated
with diffeomorphisms in GR has profound consequences for: - The physical
interpretation of spacetime points. - The definition of observables. - The
distinction between physical and gauge degrees of freedom. - The initial-value
formulation of GR. - The quantization of gravity. See the detailed review by
James Anderson and, especially, the review by John Norton. Springer Nature
Link+1 9. An important qualification concerning Kretschmann himself There is a
historical subtlety worth mentioning. The familiar textbook formulation of
Kretschmann's objection is: "Any physical theory can be formulated in generally
covariant form." This is a good summary of the central objection, but it is not
a complete account of his original argument. Kretschmann also engaged with
Einstein's point-coincidence argument and investigated the possibility of
identifying spacetime points through invariant geometrical quantities. His
argument was therefore not merely a complaint about tensor notation. The
historical discussion is quite interesting, and Norton's 1993 review is a
particularly valuable reference. IOPscience+1 10. My main conclusion: what
should we actually demand of a physical theory? I would formulate the lesson as
follows. There are three increasingly substantive requirements: 1. Covariant
formulation: The equations can be expressed independently of the choice of
coordinates. 2. Physical symmetry: The theory possesses transformations that
preserve its physical laws and structures. 3. Relational or
background-independent formulation: The theory does not rely on prescribed
geometrical structures that determine physical spacetime independently of its
dynamical contents. These requirements are not equivalent. In particular, the
third is not a simple consequence of the first or second. A theory can be
generally covariant while retaining absolute background structures. And a
theory can have physical symmetries without being generally covariant in its
original formulation. 11. A connection with Hamiltonian mechanics There is a
particularly interesting analogy with constrained Hamiltonian mechanics.
Consider a Hamiltonian system with first-class constraints: \[ \phi_a(q,p)=0.
\] The constraints generate transformations in phase space. Some
transformations connect different mathematical representations of the same
physical state. The physical state space is then obtained by quotienting the
constraint surface by the gauge transformations. Schematically, \[ \mathcal
P_{\rm phys} = \frac{\{\phi_a=0\}}{\text{gauge transformations}}. \] In GR, the
diffeomorphism freedom is reflected in the Hamiltonian and momentum
constraints. Thus, the distinction between coordinate freedom and physical
degrees of freedom becomes a question about the structure of the constrained
dynamical system. This is not accidental. General relativity is a particularly
rich example of a generally covariant gauge theory. References for further
exploration I would suggest the following reading order: 1. J. D. Norton
(1993), General covariance and the foundations of general relativity: eight
decades of dispute, Reports on Progress in Physics 56, 791. Article. The best
historical and conceptual overview. 2. J. D. Norton (2003), General covariance,
gauge theories, and the Kretschmann objection, in Symmetries in Physics:
Philosophical Reflections, pp. 110–123. Chapter. Particularly important for
reconciling Kretschmann's objection with the modern gauge interpretation. 3.
The Hole Argument and Some Physical and Philosophical Implications, Living
Reviews in Relativity. Review. A detailed discussion of diffeomorphisms,
spacetime individuation and the hole argument. One question I think deserves a
separate investigation: Is the essential physical content of general relativity
its general covariance, or is it the dynamical character of spacetime geometry?
I would argue that the second is much closer to the heart of the matter. And
this leads to a deeper question relevant to your interests in classical
mechanics and the foundations of quantum mechanics: Can a theory be
background-independent without being generally covariant, and can general
covariance be formulated without introducing gauge redundancy? Those questions
help separate the physics of spacetime from the mathematics used to describe
it. continue 12. Beyond Kretschmann: is general covariance physically empty? We
can now push the discussion further. I think the most interesting lesson of
Kretschmann's observation is not that general covariance is unimportant, but
that we must identify exactly where the physical content of a theory resides.
Let us investigate this through three questions: 1. Can a theory be
background-independent without being generally covariant? 2. Can a theory be
generally covariant without possessing genuine gauge freedom? 3. What is the
relationship between general covariance, gauge symmetry, and physical degrees
of freedom? These questions are closely related, but their answers are not
identical. 12.1 Background independence without general covariance First, we
need to clarify what we mean by background independence. A theory is called
background-independent, in a commonly used sense, when the geometrical
structures governing its dynamics are not prescribed in advance but are
themselves dynamical. For example, in GR, \[ S[g,\psi]= \frac{c^3}{16\pi G}
\int d^4x\,\sqrt{-g}\,R +S_{\mathrm m}[g,\psi]. \] Here both matter fields
\(\psi\) and the metric \(g_{\mu\nu}\) are dynamical. Compare this with a
matter theory on a prescribed spacetime: \[ S[\psi;\eta]= \int d^4x\,\mathcal
L(\psi,\partial\psi;\eta). \] The semicolon emphasizes that \(\eta_{\mu\nu}\)
is not varied. The distinction is physical, not merely notational. However,
does background independence necessarily require general covariance? Not as a
matter of pure logic. One could formulate a theory using a preferred coordinate
system, with dynamical geometrical variables and equations that refer
explicitly to those coordinates. Such a theory could lack manifest general
covariance while having no prescribed spacetime metric. Of course, we would
then have to investigate whether the preferred coordinates represent genuine
physical structure or merely an inconvenient choice of description. This
distinction is important: background independence is a statement about the
dynamical status of structures, whereas covariance is a statement about the
form of equations. 12.2 Can general covariance exist without genuine gauge
freedom? Yes. This is perhaps the most direct way to understand the force of
Kretschmann's objection. Consider a scalar field on Minkowski spacetime: \[
    S[\phi;\eta]= \frac12\int d^4x\,\sqrt{-\eta} \left(
    \eta^{\mu\nu}\partial_\mu\phi\partial_\nu\phi -m^2\phi^2 \right). \] We can
formulate this theory in arbitrary coordinates, replacing the constant
Minkowski metric by \(g_{\mu\nu}(x)\), subject to the flatness condition. The
resulting equations are generally covariant. But there is no new gravitational
gauge field. The metric is still prescribed. There is an instructive
distinction here. If we perform an active diffeomorphism on both the scalar
field and the background metric, we obtain another representation of the same
geometrical situation. But if we transform the scalar field while holding the
background metric fixed, we generally obtain a physically different
configuration, not a gauge-equivalent one. Thus, covariance of the complete
geometrical description does not automatically establish that the matter
configurations related by arbitrary diffeomorphisms are gauge-equivalent. This
is one of the central points in Norton's analysis of the Kretschmann objection.
sites.pitt.edu+1 12.3 The role of Noether's second theorem Here is a more
technical way to distinguish ordinary symmetries from gauge symmetries.
Noether's first theorem concerns continuous transformations depending on
finitely many constant parameters. For example, translation invariance leads to
energy-momentum conservation. Noether's second theorem concerns transformations
depending on arbitrary functions of spacetime. Schematically, suppose the
fields \(\Psi^A\) transform as \[ \delta\Psi^A = R^A_{\alpha}\epsilon^\alpha +
R^{A\mu}_{\alpha}\partial_\mu\epsilon^\alpha +\cdots, \] where the
\(\epsilon^\alpha(x)\) are arbitrary functions. If the action is invariant
under these transformations, the Euler–Lagrange expressions satisfy
differential identities: \[ \boxed{ R^{A}_{\alpha}E_A
-\partial_\mu(R^{A\mu}_{\alpha}E_A) +\cdots=0. } \] Here \[ E_A=\frac{\delta
S}{\delta\Psi^A}. \] These are Noether identities. They express dependencies
among the field equations. In electromagnetism, for example, \[
    A_\mu\longrightarrow A_\mu+\partial_\mu\chi \] is a gauge transformation.
The field equations are \[ \partial_\mu F^{\mu\nu}=j^\nu, \] where \[
    F_{\mu\nu}=\partial_\mu A_\nu-\partial_\nu A_\mu. \] Taking the divergence
gives \[ \partial_\nu j^\nu=0. \] The conservation of charge is therefore
related to the differential identity arising from gauge invariance. In general
relativity, the corresponding identity is the contracted Bianchi identity: \[
    \boxed{\nabla_\mu G^{\mu\nu}=0.} \] Einstein's equations consequently imply
\[ \nabla_\mu T^{\mu\nu}=0. \] This is a profound relationship between
spacetime geometry, gauge freedom and the structure of the equations of motion.
Notice, however, that the Bianchi identity is a geometrical identity. It does
not depend on imposing Einstein's equations. This distinction between
identities and equations of motion is essential. 13. General relativity as a
constrained dynamical system Now let us connect this with Hamiltonian
mechanics, where the distinction becomes particularly concrete. In the ADM
formulation, spacetime is foliated into spatial hypersurfaces: \[ ds^2=N^2dt^2
-h_{ij}(dx^i+N^i dt)(dx^j+N^jdt), \] using the \(+---\) spacetime signature.
Here: - \(h_{ij}\) is the spatial metric; - \(N\) is the lapse; - \(N^i\) is
the shift. The canonical variables are \[ (h_{ij},\pi^{ij}). \] The lapse and
shift act as Lagrange multipliers. Their variation gives the constraints: \[
    \mathcal H_\perp=0, \qquad \mathcal H_i=0. \] The Hamiltonian is \[ H_{\rm
ADM} = \int d^3x \left( N\mathcal H_\perp+N^i\mathcal H_i \right)
+H_{\partial}. \] The constraints generate deformations of the spatial
hypersurface, corresponding to spacetime diffeomorphisms, subject to important
qualifications concerning boundaries and the precise generators. The constraint
algebra is \[ \{H[\vec N],H[\vec M]\} =H[[\vec N,\vec M]], \] \[ \{H[\vec
N],H_\perp[M]\} =H_\perp[\mathcal L_{\vec N}M], \] \[ \{H_\perp[N],H_\perp[M]\}
= H\left[ h^{ij}(N\partial_jM-M\partial_jN) \right]. \] The last relation is
especially interesting. It is not simply the Lie algebra of a fixed spacetime
symmetry group. Its structure functions depend on the dynamical spatial metric.
This is called the hypersurface-deformation algebra. It reflects the geometry
of embedding one spatial hypersurface into spacetime. Counting physical degrees
of freedom The spatial metric has six independent components at each point, and
so does its conjugate momentum. Thus there are twelve phase-space variables.
There are four first-class constraints: \[ \mathcal H_\perp=0,\qquad \mathcal
H_i=0. \] Each independent first-class constraint removes two phase-space
dimensions: one through the constraint and one through the associated gauge
freedom. Consequently, \[ 12-2(4)=4 \] physical phase-space dimensions per
spatial point, or \[ \boxed{2\text{ configuration-space degrees of freedom}.}
\] These are the two propagating gravitational-wave polarizations. This is a
much more substantive statement than saying that Einstein's equations are
generally covariant. It tells us how many independent gravitational degrees of
freedom the theory possesses. 14. An important subtlety: not every
diffeomorphism is physically gauge There is a qualification to the usual
statement that diffeomorphisms in GR are gauge transformations. Consider an
asymptotically flat spacetime. We usually impose boundary conditions at spatial
infinity. Some diffeomorphisms preserve these boundary conditions, but they do
not necessarily represent pure gauge transformations. For example, asymptotic
time translations can carry a nonzero ADM energy. Similarly, asymptotic spatial
translations can carry momentum. The distinction is: - Proper gauge
transformations: transformations that leave the physical state unchanged. -
Large or asymptotic symmetries: transformations that may act nontrivially on
physical states and carry conserved charges. Thus, even within general
relativity, the statement that all diffeomorphisms are gauge requires
qualification. The boundary conditions and the definition of the physical phase
space matter. This becomes especially important in gravitational radiation,
black-hole mechanics and the asymptotic symmetry group of spacetime. 15.
Comparison with electromagnetism I think comparing electromagnetism and GR is
particularly illuminating. Feature	Electromagnetism	General relativity
Gauge variable	\(A_\mu\)	\(g_{\mu\nu}\) and spacetime description Gauge
transformation	\(A_\mu\to A_\mu+\partial_\mu\chi\)	Diffeomorphisms
Gauge-invariant field	\(F_{\mu\nu}\)	Curvature and relational geometrical
quantities Background metric	Usually prescribed	Dynamical First-class
constraints	Gauss constraint	Hamiltonian and momentum constraints Local
propagating DOF	2	2 There is an important difference, however.
Electromagnetism has an internal gauge group \(U(1)\), whereas GR's
diffeomorphism freedom concerns spacetime itself. The two theories therefore
have different geometrical interpretations, even though their constrained
Hamiltonian structures have important similarities. This also explains why the
general covariance of GR should not simply be identified with the local gauge
invariance of electromagnetism. 16. What does this mean for the physical
interpretation of spacetime? We can now return to the original question.
Kretschmann showed that the mathematical form of general covariance cannot by
itself establish the physical content of GR. But the subsequent development of
the theory revealed something much deeper: The metric is not merely a
geometrical background on which physical processes occur. It is itself a
physical field whose dynamics participate in determining spacetime geometry.
Moreover, the diffeomorphism freedom means that individual manifold points
cannot be assigned independent physical identities. Physical predictions must
ultimately be formulated in terms of invariant or relational quantities. For
example, rather than asking for the value of a field at a bare manifold point
\(x\), one might ask for its value where some physical clock field takes the
value \(\tau\), and some physical spatial reference fields take specified
values. This is one route towards relational observables. There is a
difficulty: constructing sufficiently general and useful relational observables
is not straightforward. But it is precisely here that the physical content of
general relativity goes beyond Kretschmann's formal observation. 17. My
synthesis I would summarize the entire discussion in the following hierarchy:
Level 1 — Coordinate independence The physical predictions do not depend on the
coordinate system used to describe them. Level 2 — General covariance The
equations admit a formulation valid in arbitrary smooth coordinates. Level 3 —
Gauge freedom Different mathematical configurations can represent the same
physical state. Level 4 — Background independence The geometrical structures
defining spacetime are not prescribed independently of the dynamics. Level 5 —
Relational physical content Physical observables are defined through relations
among dynamical fields and physical reference systems. These levels are
related, but they are not logically equivalent. In particular, Kretschmann's
observation concerns Level 2, whereas much of the distinctive physics of GR
resides in Levels 3–5. Further reading Three particularly useful references: 1.
J. D. Norton, General Covariance, Gauge Theories and the Kretschmann Objection.
This directly addresses the apparent contradiction between Kretschmann's
objection and the gauge interpretation. 2. J. Stachel, The Hole Argument and
Some Physical and Philosophical Implications. Particularly useful for the
relationship between diffeomorphisms, spacetime individuation and background
independence. 3. Stanford Encyclopedia of Philosophy: Symmetry and Symmetry
Breaking. Its discussion of general covariance and absolute objects is useful
for placing the issue in the broader theory of symmetry. sites.pitt.edu+2 18. A
question I would pursue next There is one issue that I think deserves
particular attention in light of your interest in classical mechanics,
reduction and constraints. Is general covariance fundamentally a symmetry
principle, or is it fundamentally a statement about the redundancy of the
dynamical description? These are not quite the same. In ordinary mechanics,
symmetry reduction starts with a physical symmetry group acting on phase space.
We then construct a reduced phase space. In GR, the diffeomorphism group is
involved in the very definition of physical configurations, and the constraint
algebra is not an ordinary Lie algebra with constant structure coefficients.
This raises a deeper question: Can we understand general relativity as a
particular kind of dynamical reduction, in which spacetime geometry is
constructed only after eliminating the redundant variables? That would connect
Kretschmann's observation directly with symplectic reduction, Dirac
constraints, and the physical meaning of gauge symmetry.
